\naslov{Toplotna enačba}

Naj bo $U \subseteq \R^n$ omejena odprta množica z gladkim robom.
Obravnavamo Cauchyjev problem za toplotno enačbo
\[
  u_t - \Delta u = 0
\]
na $U \times (0,\infty)$.
Poleg tega zahtevamo $u=0$ na $\partial U \times (0,\infty)$ in $u(x,0) = g(x)$
na $U \times \{0\}$.
Iskali bomo rešitve oblike $u(x,t) = v(t) w(x)$.
Izračunamo $u_t - \Delta u = w(x) v'(t) - v(t) \Delta w(x) = 0$, iz česar dobimo
\[
  \frac{\Delta w(x)}{w(x)} = \frac{v'(t)}{v(t)}.
\]
Leva stran te enačbe je neodvisna od $t$, desna pa neodvisna od $x$, torej
morata biti konstantni.
Označimo to konstanto z $-\lambda$.
Iskali bomo funkcije $w$, ki rešijo Cauchyjevo nalogo $- \Delta w = \lambda w$
in $w = 0$ na $\partial U$.
Potem bo funkcija $u(x,t) = d e^{-\lambda t} w(x)$ rešila enačbo $u_t - \Delta u
= 0$ na $U \times (0,\infty)$ s pogojem $u = 0$ na $\partial U$.
Pri $t = 0$ velja $u(x,0) = d w(x) = g(x)$, torej je rešitev oblike $u(x,t) =
g(x) e^{-\lambda t}$, a le v primeru, da je $g$ oblike $d w(x)$, kjer je $d \in
\R$ in $w$ reši omenjen problem.
Ker je enačba linearna, je vsota rešitev take oblike tudi rešitev.
Poiskali bomo števno družino rešitev $w_j, \lambda_j, d_j$, da bo
\[
  u(x,t) = \sum_{j=1}^\infty d_j e^{-\lambda_j t} w_j(x).
\]

\vprasanje{Izpelji splošno rešitev za toplotno enačbo.}

Obravnavajmo primer $n=1$.
Če je $\lambda \le 0$, iz robnih pogojev izpeljemo, da mora biti $w$ trivialen,
kar za nas ni zanimivo.
V primeru $\lambda > 0$ dobimo
\[
  w(x) = A \cos(\sqrt{\lambda} x) + B \sin(\sqrt{\lambda} x).
\]
Iz robnih pogojev dobimo
\[
  \begin{bmatrix}
	\cos(\sqrt{\lambda} a) & \sin(\sqrt{\lambda} a) \\
	\cos(\sqrt{\lambda} b) & \sin(\sqrt{\lambda} b)
  \end{bmatrix}
  \cdot
  \begin{bmatrix}
	A \\ B
  \end{bmatrix}
  = 0.
\]
Dobimo dodatno zahtevno, da je ta matrika nesingularna.
Sistem se prevede na $\sin(\sqrt{\lambda} (b-a)) = 0$, torej mora veljati
\[
  \lambda = \lambda_k = \left( \frac{k \pi}{b-a} \right)^2.
\]
Upoštevamo $a = 0$ in $b = L$, da dobimo $\lambda_k = (k \pi / L)^2$.
Iz $w(a) = w(b) = 0$ potem sledi $A = 0$, in so rešitve oblike
\[
  w_k(x) = B_k \sin(\frac{k \pi x}{L}).
\]
Nastavek sedaj deluje za
\[
  g = \sum_{k=1}^\infty d_k w_k,
\]
in iščemo koeficiente $B_k$, da bo ta enakost veljala.

\vprasanje{Izpelji rešitev Cauchyjevega problema za enodimenzionalno toplotno
  enačbo.}

% LocalWords:  Cauchyjevega
